\chapter{Introduction to Combinational Logic Circuits}
\label{chap:combinational-logic}

\section{Introduction \& Design Principles}
Combinational logic circuits are the fundamental building blocks of digital processing architectures. In a combinational circuit, the outputs at any instant of time depend **strictly on the present combination of input variables**, with no influence from past inputs or stored states.

\begin{definition}[Combinational Logic Circuit]
A network of interconnected digital logic gates whose outputs at any given time are determined solely by the instantaneous values of its inputs, without feedback paths or internal memory elements.
\end{definition}

\noindent **The Standard Combinational Design Flow:**
\begin{center}
\begin{tikzpicture}[node distance=1.0cm, auto, >=Stealth, transform shape, scale=0.75]
    \node (p1) [draw=brandnavy, fill=boxnavybg, rounded corners, inner sep=5pt, font=\small\bfseries] {1. Specification};
    \node (p2) [draw=brandblue, fill=boxskybg, rounded corners, inner sep=5pt, right=0.6cm of p1, font=\small\bfseries] {2. Truth Table};
    \node (p3) [draw=brandamber, fill=boxgoldbg, rounded corners, inner sep=5pt, right=0.6cm of p2, font=\small\bfseries] {3. K-Map Minimization};
    \node (p4) [draw=brandteal, fill=boxtealbg, rounded corners, inner sep=5pt, right=0.6cm of p3, font=\small\bfseries] {4. Circuit Synthesis};
    
    \draw[->, line width=1.3pt, brandnavy] (p1) -- (p2);
    \draw[->, line width=1.3pt, brandblue] (p2) -- (p3);
    \draw[->, line width=1.3pt, brandamber] (p3) -- (p4);
\end{tikzpicture}
\end{center}

---

\section{Arithmetic Adders}

\subsection{Half Adder}
A Half Adder is an arithmetic circuit that adds two single-bit binary inputs ($A$ and $B$) and generates two output bits: **Sum ($S$)** and **Carry ($C$)**.

\begin{center}
\begin{circuitikz}[scale=0.85, transform shape]
    % XOR Gate for Sum
    \node (xor1) at (2, 2.2) [xor gate US, draw, logic gate inputs=nn] {};
    \draw (xor1.input 1) -- ++(-1.5,0) node[left] {$A$};
    \draw (xor1.input 2) -- ++(-1.5,0) node[left] {$B$};
    \draw (xor1.output) -- ++(1.0,0) node[right, font=\bfseries] {$S = A \oplus B$};

    % AND Gate for Carry
    \node (and1) at (2, 0.4) [and gate US, draw, logic gate inputs=nn] {};
    \draw (xor1.input 1) ++(-0.8,0) node[circ] {} |- (and1.input 1);
    \draw (xor1.input 2) ++(-0.5,0) node[circ] {} |- (and1.input 2);
    \draw (and1.output) -- ++(1.0,0) node[right, font=\bfseries] {$C = A \cdot B$};
\end{circuitikz}
\end{center}

\begin{table}[htbp]
\centering
\small
\begin{tabular}{cc|cc|l}
\toprule
\textbf{Input $A$} & \textbf{Input $B$} & \textbf{Carry ($C$)} & \textbf{Sum ($S$)} & \textbf{Arithmetic Addition} \\
\midrule
0 & 0 & 0 & 0 & $0 + 0 = 00_2$ \\
0 & 1 & 0 & 1 & $0 + 1 = 01_2$ \\
1 & 0 & 0 & 1 & $1 + 0 = 01_2$ \\
1 & 1 & 1 & 0 & $1 + 1 = 10_2$ \\
\bottomrule
\end{tabular}
\caption{Truth table for a Binary Half Adder.}
\label{tab:half-adder}
\end{table}

\begin{keyequation}[Half Adder Boolean Expressions]
\begin{equation}
    S = A \oplus B = \overline{A}B + A\overline{B}, \qquad C = A \cdot B
\end{equation}
\end{keyequation}

\subsection{Full Adder}
A Half Adder cannot accept a carry generated from a lower significant column. A **Full Adder** adds three 1-bit inputs: $A$, $B$, and an incoming carry $C_{in}$, producing Sum ($S$) and Carry-Out ($C_{out}$).

\begin{table}[htbp]
\centering
\small
\begin{tabular}{ccc|cc|l}
\toprule
\textbf{Input $A$} & \textbf{Input $B$} & \textbf{Input $C_{in}$} & \textbf{Carry Out ($C_{out}$)} & \textbf{Sum ($S$)} & \textbf{Decimal Value} \\
\midrule
0 & 0 & 0 & 0 & 0 & $0+0+0 = 0$ \\
0 & 0 & 1 & 0 & 1 & $0+0+1 = 1$ \\
0 & 1 & 0 & 0 & 1 & $0+1+0 = 1$ \\
0 & 1 & 1 & 1 & 0 & $0+1+1 = 2$ \\
1 & 0 & 0 & 0 & 1 & $1+0+0 = 1$ \\
1 & 0 & 1 & 1 & 0 & $1+0+1 = 2$ \\
1 & 1 & 0 & 1 & 0 & $1+1+0 = 2$ \\
1 & 1 & 1 & 1 & 1 & $1+1+1 = 3$ \\
\bottomrule
\end{tabular}
\caption{Truth table for a Binary Full Adder.}
\label{tab:full-adder}
\end{table}

\begin{keyequation}[Full Adder Boolean Equations]
\begin{align}
    S &= A \oplus B \oplus C_{in} \\
    C_{out} &= AB + BC_{in} + AC_{in} = AB + C_{in}(A \oplus B)
\end{align}
\end{keyequation}

\begin{figure}[htbp]
\centering
\begin{tikzpicture}[scale=0.85, transform shape]
    % Half Adder 1
    \draw[fill=boxskybg, draw=boxskyborder, line width=1.2pt, rounded corners=4pt] (0,1) rectangle (3.2, 3.5);
    \node at (1.6, 3.1) [font=\small\bfseries\color{brandblue}] {Half Adder 1};
    \draw (0, 2.5) -- ++(-0.8,0) node[left] {$A$};
    \draw (0, 1.7) -- ++(-0.8,0) node[left] {$B$};
    
    % Half Adder 2
    \draw[fill=boxskybg, draw=boxskyborder, line width=1.2pt, rounded corners=4pt] (4.5, 0.5) rectangle (7.7, 3.0);
    \node at (6.1, 2.6) [font=\small\bfseries\color{brandblue}] {Half Adder 2};
    \draw (4.5, 1.2) -- ++(-5.3,0) node[left] {$C_{in}$};
    
    % Connections
    \draw (3.2, 2.2) -- (4.5, 2.2) node[midway, above, font=\scriptsize] {$A \oplus B$};
    \draw (7.7, 2.0) -- ++(1.5,0) node[right, font=\bfseries] {$S = A \oplus B \oplus C_{in}$};
    
    % OR gate for Carry
    \node (or1) at (9.5, 0.5) [or gate US, draw, logic gate inputs=nn] {};
    \draw (3.2, 1.3) -- ++(0.5,0) |- (or1.input 1) node[pos=0.2, above, font=\tiny] {$C_1{=}AB$};
    \draw (7.7, 0.8) -- ++(0.5,0) |- (or1.input 2) node[pos=0.2, below, font=\tiny] {$C_2{=}C_{in}(A\oplus B)$};
    \draw (or1.output) -- ++(0.8,0) node[right, font=\bfseries] {$C_{out}$};
\end{tikzpicture}
\caption{Full Adder synthesized using Two Half Adders and One OR Gate.}
\label{fig:full-adder-two-ha}
\end{figure}

---

\section{Arithmetic Subtractors}

\subsection{Half Subtractor}
Subtracts a 1-bit subtrahend ($B$) from a 1-bit minuend ($A$), generating **Difference ($D$)** and **Borrow-Out ($B_{out}$)**.

\begin{table}[htbp]
\centering
\small
\begin{tabularx}{\textwidth}{cc|cc|X}
\toprule
\textbf{Minuend ($A$)} & \textbf{Subtrahend ($B$)} & \textbf{Borrow ($B_{out}$)} & \textbf{Difference ($D$)} & \textbf{Arithmetic Operation ($A - B$)} \\
\midrule
0 & 0 & 0 & 0 & $0 - 0 = 0$ \\
0 & 1 & 1 & 1 & $0 - 1 = -1 \implies \text{Borrow } 1 \rightarrow \text{Diff } 1$ \\
1 & 0 & 0 & 1 & $1 - 0 = 1$ \\
1 & 1 & 0 & 0 & $1 - 1 = 0$ \\
\bottomrule
\end{tabularx}
\caption{Truth table for a Binary Half Subtractor.}
\label{tab:half-subtractor}
\end{table}

\begin{keyequation}[Half Subtractor Boolean Equations]
\begin{equation}
    D = A \oplus B, \qquad B_{out} = \overline{A} \cdot B
\end{equation}
\end{keyequation}

\subsection{Full Subtractor}
Performs binary subtraction involving three bits: minuend ($A$), subtrahend ($B$), and borrow-in ($B_{in}$).

\begin{keyequation}[Full Subtractor Boolean Equations]
\begin{align}
    D &= A \oplus B \oplus B_{in} \\
    B_{out} &= \overline{A}B + \overline{A}B_{in} + BB_{in} = \overline{A}B + B_{in}\overline{(A \oplus B)}
\end{align}
\end{keyequation}

\noindent A Full Subtractor can be realized by cascading **two Half Subtractors and one OR gate**.

---

\section{Multiplexers (MUX / Data Selectors)}

\subsection{Theory of Operation}
A Multiplexer (MUX) is a combinational circuit that selects binary information from one of $2^n$ input lines and directs it to a single output line based on the binary state of $n$ selection lines.

\begin{keyequation}[Multiplexer Fundamental Dimensioning]
\begin{equation}
    N_{\text{inputs}} = 2^n \iff n = \log_2(N_{\text{inputs}})
\end{equation}
Where $N_{\text{inputs}}$ is the number of data input lines and $n$ is the number of select lines.
\end{keyequation}

\subsection{4:1 Multiplexer}
\begin{itemize}
    \item \textbf{Data Inputs:} $I_0, I_1, I_2, I_3$ ($2^2 = 4 \implies n=2$)
    \item \textbf{Select Lines:} $S_1, S_0$ ($S_1$ is MSB, $S_0$ is LSB)
    \item \textbf{Output:} $Y$
\end{itemize}

\begin{table}[htbp]
\centering
\small
\begin{tabular}{cc|c|l}
\toprule
\textbf{Select $S_1$} & \textbf{Select $S_0$} & \textbf{Output $Y$} & \textbf{Selected Channel} \\
\midrule
0 & 0 & $I_0$ & Channel 0 Selected \\
0 & 1 & $I_1$ & Channel 1 Selected \\
1 & 0 & $I_2$ & Channel 2 Selected \\
1 & 1 & $I_3$ & Channel 3 Selected \\
\bottomrule
\end{tabular}
\caption{Function Table for a 4:1 Multiplexer.}
\label{tab:4-to-1-mux}
\end{table}

\begin{keyequation}[4:1 Multiplexer Output Equation]
\begin{equation}
    Y = \overline{S_1}\,\overline{S_0}I_0 + \overline{S_1}S_0I_1 + S_1\overline{S_0}I_2 + S_1S_0I_3
\end{equation}
\end{keyequation}

\begin{figure}[htbp]
\centering
\begin{circuitikz}[scale=0.85, transform shape]
    % 4 AND Gates
    \node (and0) at (2,3.5) [and gate US, draw, logic gate inputs=nnn, scale=0.8] {};
    \node (and1) at (2,2.3) [and gate US, draw, logic gate inputs=nnn, scale=0.8] {};
    \node (and2) at (2,1.1) [and gate US, draw, logic gate inputs=nnn, scale=0.8] {};
    \node (and3) at (2,-0.1) [and gate US, draw, logic gate inputs=nnn, scale=0.8] {};
    
    % OR Gate
    \node (or1) at (5,1.7) [or gate US, draw, logic gate inputs=nnnn, scale=1.1] {};
    
    \draw (and0.output) -- ++(0.8,0) |- (or1.input 1);
    \draw (and1.output) -- ++(0.4,0) |- (or1.input 2);
    \draw (and2.output) -- ++(0.4,0) |- (or1.input 3);
    \draw (and3.output) -- ++(0.8,0) |- (or1.input 4);
    
    \draw (or1.output) -- ++(1.0,0) node[right, font=\bfseries] {$Y$};
    
    \draw (and0.input 1) -- ++(-2.5,0) node[left] {$I_0$};
    \draw (and1.input 1) -- ++(-2.5,0) node[left] {$I_1$};
    \draw (and2.input 1) -- ++(-2.5,0) node[left] {$I_2$};
    \draw (and3.input 1) -- ++(-2.5,0) node[left] {$I_3$};
\end{circuitikz}
\caption{Gate-level logic implementation of a 4-to-1 Multiplexer.}
\label{fig:4-to-1-mux-circuit}
\end{figure}

\begin{example}[Boolean Function Synthesis using an 8:1 Multiplexer]
Implement the 3-variable logic function $F(A,B,C) = \sum m(1, 3, 4, 6, 7)$ using an 8:1 Multiplexer.
\tcblower
\textbf{Step 1: Determine Dimensions}
Number of variables $n = 3 \implies 2^3 = 8$ data inputs ($I_0$ to $I_7$) and $3$ select lines ($S_2 = A, S_1 = B, S_0 = C$).

\textbf{Step 2: Map Minterms to Multiplexer Data Inputs}
\begin{itemize}
    \item Connect minterm indices present in $F$ ($1, 3, 4, 6, 7$) to logic HIGH ($\mathbf{1} / V_{CC}$).
    \item Connect remaining minterm indices ($0, 2, 5$) to logic LOW ($\mathbf{0} / \text{GND}$).
\end{itemize}

\textbf{Step 3: Multiplexer Connections}
$$\mathbf{I_1 = I_3 = I_4 = I_6 = I_7 = 1 \text{ (5V)}, \qquad I_0 = I_2 = I_5 = 0 \text{ (GND)}}$$
$$\mathbf{S_2 = A, \quad S_1 = B, \quad S_0 = C}$$
\end{example}

---

\section{De-Multiplexers (DEMUX / Data Distributors)}

\subsection{Theory of Operation}
A De-multiplexer receives data on a single input line and steers it to one of $2^n$ output lines depending on the values of $n$ selection lines.

\begin{keyequation}[1:4 De-multiplexer Equations]
\begin{align}
    Y_0 &= \overline{S_1}\,\overline{S_0} \cdot D_{in} \\
    Y_1 &= \overline{S_1}S_0 \cdot D_{in} \\
    Y_2 &= S_1\overline{S_0} \cdot D_{in} \\
    Y_3 &= S_1S_0 \cdot D_{in}
\end{align}
\end{keyequation}

---

\section{Decoders and Encoders}

\subsection{Binary Decoders (2-to-4 and 3-to-8 Lines)}
A Decoder is a combinational circuit that converts an $n$-bit binary code into up to $2^n$ unique output lines. 

\begin{itemize}
    \item **Active-High Outputs:** The selected line outputs logic $\mathbf{1}$; all non-selected lines output logic $\mathbf{0}$.
    \item **Active-Low Outputs (e.g., 74138 Decoder):** The selected line outputs logic $\mathbf{0}$; all non-selected lines output logic $\mathbf{1}$.
\end{itemize}

\begin{table}[htbp]
\centering
\small
\begin{tabular}{c|cc|cccc}
\toprule
\textbf{Enable ($E$)} & \textbf{Input $A$} & \textbf{Input $B$} & \textbf{$Y_3$} & \textbf{$Y_2$} & \textbf{$Y_1$} & \textbf{$Y_0$} \\
\midrule
0 & X & X & 0 & 0 & 0 & 0 \\
1 & 0 & 0 & 0 & 0 & 0 & 1 \\
1 & 0 & 1 & 0 & 0 & 1 & 0 \\
1 & 1 & 0 & 0 & 1 & 0 & 0 \\
1 & 1 & 1 & 1 & 0 & 0 & 0 \\
\bottomrule
\end{tabular}
\caption{Truth table for a 2-to-4 Line Decoder with Active-High outputs.}
\label{tab:2-to-4-decoder}
\end{table}

\subsection{Binary Encoders \& Priority Encoders}
An Encoder performs the inverse operation of a decoder: it accepts $2^n$ (or fewer) input lines and generates an $n$-bit binary code.

\subsubsection{8-to-3 Octal-to-Binary Encoder}
\begin{itemize}
    \item \textbf{Inputs:} $D_0, D_1, D_2, D_3, D_4, D_5, D_6, D_7$ (only one input is active at any time).
    \item \textbf{Outputs:} $X, Y, Z$ ($X$ is MSB, $Z$ is LSB).
\end{itemize}

\begin{keyequation}[8-to-3 Encoder Output Equations]
\begin{align}
    Z &= D_1 + D_3 + D_5 + D_7 \\
    Y &= D_2 + D_3 + D_6 + D_7 \\
    X &= D_4 + D_5 + D_6 + D_7
\end{align}
\end{keyequation}

\begin{commonmistake}[Priority Encoder Requirement]
Standard encoders produce invalid outputs when more than one input is active simultaneously. A **Priority Encoder** establishes an explicit input priority (e.g., $D_7 > D_6 > \dots > D_0$). If both $D_2$ and $D_6$ are active, the priority encoder exclusively outputs code $110_2$ (corresponding to $D_6$). It also provides a **Valid Output ($V$)** flag to differentiate input $D_0 = 1$ from an all-zero inactive condition.
\end{commonmistake}

---

\section{Digital Magnitude Comparators}

\subsection{1-Bit Magnitude Comparator}
Compares two 1-bit inputs $A$ and $B$, generating three mutually exclusive outputs:
\begin{keyequation}[1-Bit Magnitude Comparator Equations]
\begin{align}
    (A > B) &= A \cdot \overline{B} \\
    (A = B) &= A \odot B = \overline{A \oplus B} = AB + \overline{A}\,\overline{B} \\
    (A < B) &= \overline{A} \cdot B
\end{align}
\end{keyequation}

\subsection{2-Bit Magnitude Comparator}
Compares two 2-bit words: $A = (A_1 A_0)_2$ and $B = (B_1 B_0)_2$.

\begin{figure}[htbp]
\centering
\begin{circuitikz}[scale=0.85, transform shape]
    % Comparator Box
    \draw[fill=boxnavybg, draw=brandnavy, line width=1.5pt, rounded corners=6pt] (0,0) rectangle (5,4);
    \node at (2.5, 3.5) [font=\bfseries\sffamily\color{brandnavy}] {2-Bit Comparator};
    
    % Inputs
    \draw (-1.5, 2.8) node[left] {$A_1 A_0$} -- (0, 2.8);
    \draw (-1.5, 1.2) node[left] {$B_1 B_0$} -- (0, 1.2);
    
    % Outputs
    \draw (5, 3.0) -- ++(1.5,0) node[right, font=\bfseries] {$A > B$};
    \draw (5, 2.0) -- ++(1.5,0) node[right, font=\bfseries] {$A = B$};
    \draw (5, 1.0) -- ++(1.5,0) node[right, font=\bfseries] {$A < B$};
\end{circuitikz}
\caption{Block representation of a 2-bit digital magnitude comparator.}
\label{fig:2-bit-comparator}
\end{figure}

\begin{keyequation}[2-Bit Magnitude Comparator Equations]
Let equality flags for individual bit positions be:
\begin{equation}
    x_1 = \overline{A_1 \oplus B_1} = A_1B_1 + \overline{A_1}\,\overline{B_1}, \qquad x_0 = \overline{A_0 \oplus B_0} = A_0B_0 + \overline{A_0}\,\overline{B_0}
\end{equation}
Then the comparator outputs are synthesized as:
\begin{align}
    (A = B) &= x_1 \cdot x_0 \\
    (A > B) &= A_1\overline{B_1} + x_1 (A_0\overline{B_0}) \\
    (A < B) &= \overline{A_1}B_1 + x_1 (\overline{A_0}B_0)
\end{align}
\end{keyequation}

---

\section{Summary \& Chapter Review}
\begin{quickrevision}[Unit 4 Essential Review Sheet]
\begin{itemize}
    \item \textbf{Combinational Circuit:} Memoryless; output depends strictly on instantaneous inputs.
    \item \textbf{Half Adder:} $S = A \oplus B$, $C = AB$.
    \item \textbf{Full Adder:} $S = A \oplus B \oplus C_{in}$, $C_{out} = AB + C_{in}(A \oplus B)$. Constructed with 2 Half Adders and 1 OR gate.
    \item \textbf{Half Subtractor:} $D = A \oplus B$, $B_{out} = \bar{A}B$.
    \item \textbf{Full Subtractor:} $D = A \oplus B \oplus B_{in}$, $B_{out} = \bar{A}B + B_{in}\overline{(A \oplus B)}$. Constructed with 2 Half Subtractors and 1 OR gate.
    \item \textbf{Multiplexer (MUX):} $2^n$ data inputs, $n$ select lines, 1 output. Acts as a universal logic synthesizer.
    \item \textbf{De-multiplexer (DEMUX):} 1 data input, $n$ select lines, $2^n$ outputs.
    \item \textbf{Decoder:} $n$ inputs $\rightarrow 2^n$ outputs. Active-High vs Active-Low (74138).
    \item \textbf{Encoder:} $2^n$ inputs $\rightarrow n$ outputs. Priority encoders resolve multi-input conflicts.
    \item \textbf{2-Bit Comparator:} Evaluates $A>B$, $A=B$, $A<B$ using bitwise equality terms $x_1, x_0$.
\end{itemize}
\end{quickrevision}
